\documentclass{beamer}

\usetheme{Madrid}
\usecolortheme{seagull}
\usepackage{amsmath, amssymb}
\usepackage{tikz}
\usepackage{bm}

\title[Integral Transforms in IBAMR]{Integral Transforms and Regularized Delta Functions in the Immersed Boundary Method}
\author{Laura A. Miller}
\date{}

\begin{document}

%=======================================================
\begin{frame}
\titlepage
\end{frame}
%=======================================================

\begin{frame}{Overview}
\tableofcontents
\end{frame}

%=======================================================
\section{The Immersed Boundary Framework}
%=======================================================

\begin{frame}{Eulerian--Lagrangian Coupling}
IB uses:
\begin{itemize}
    \item Eulerian description of fluid: velocity $\mathbf{u}(\mathbf{x},t)$, pressure $p(\mathbf{x},t)$
    \item Lagrangian description of structure: $\mathbf{X}(s,t)$
    \item Integral transforms to communicate between them
\end{itemize}

\[
\mathbf{f}(\mathbf{x},t)
=
\int_{\mathcal{U}} \mathbf{F}(s,t)\,\delta(\mathbf{x}-\mathbf{X}(s,t))\,ds
\]

\[
\frac{\partial \mathbf{X}}{\partial t}(s,t)
=
\int_\Omega \mathbf{u}(\mathbf{x},t)\,\delta(\mathbf{x}-\mathbf{X}(s,t))\,d\mathbf{x}
\]

Key idea: the delta kernel distributes forces to nearby Eulerian nodes and interpolates velocities back to the structure.
\end{frame}

\begin{frame}{Two Equivalent Lagrangian Notations}
\small

You will see two common ways to parameterize the structure:

\medskip
\textbf{(1) Curve/surface parameter $s$ (classical IB notation):}
\[
\mathbf{x}=\mathbf{X}(s,t),
\quad s\in\mathcal{U}.
\]

\textbf{(2) Reference coordinate $\mathbf{X}$ (continuum mechanics / IBFE notation):}
\[
\mathbf{x}=\boldsymbol{\chi}(\mathbf{X},t),
\quad \mathbf{X}\in\mathcal{U}.
\]

\medskip
Both represent a mapping from the reference/material domain $\mathcal{U}$ to physical space $\Omega$.
We will use $\boldsymbol{\chi}(\mathbf{X},t)$ when discussing stresses and deformation gradients.

\end{frame}

%=======================================================
\section{Force Spreading}
%=======================================================

\begin{frame}{Force Spreading: Lagrangian Reference Domain}

From the partitioned weak form (Griffith \& Luo 2017):

\[
\mathbf{f}(\mathbf{x},t)
=
\int_{\mathcal{U}}
\mathbf{F}(\mathbf{X},t)\,
\delta(\mathbf{x}-\boldsymbol{\chi}(\mathbf{X},t))\,d\mathbf{X}
\]

\medskip

\textbf{Definition of $\mathcal{U}$:}
\begin{itemize}
    \item $\mathcal{U}$ is the \emph{Lagrangian reference (material) domain} of the structure
    \item $\mathbf{X} \in \mathcal{U}$ labels material points
    \item $\mathcal{U}$ is fixed in time
    \item The mapping
    \[
    \boldsymbol{\chi} : \mathcal{U} \times [0,T] \rightarrow \Omega
    \]
    gives the current physical position of each material point
\end{itemize}

\end{frame}

\begin{frame}{What Is $\mathcal{U}$ Physically?}

\textbf{$\mathcal{U}$ represents the undeformed material description of the structure.}

\medskip

\begin{itemize}
    \item Integration is over \emph{material points}, not physical space
    \item Forces are computed in reference coordinates
    \item The delta function transfers those forces to physical space
\end{itemize}

\medskip

\textbf{Dimensionality of $\mathcal{U}$ depends on the structure:}
\begin{itemize}
    \item 1D: fibers, filaments, membranes in 2D
    \item 2D: shells or surfaces in 3D
    \item 2D or 3D: volumetric solids
\end{itemize}

\medskip

\textbf{Important distinction:}
\[
\mathcal{U}
\;\neq\;
\boldsymbol{\chi}(\mathcal{U},t)
\;\neq\;
\Omega
\]

\medskip

This separation is part of why no Jacobian appears in the force-spreading integral.

\end{frame}


%=======================================================
\section{Where Did the Jacobian Go?}
%=======================================================

\begin{frame}{Strong Form: Elastic Forces in IB}

In the immersed boundary-finite element (IBFE) framework, elasticity is described in
\emph{Lagrangian (material) coordinates}.

\medskip

\textbf{Lagrangian variables:}
\begin{itemize}
    \item $\mathbf{X} \in \mathcal{U}$: material (reference) coordinate
    \item $\boldsymbol{\chi}(\mathbf{X},t)$: current physical position
    \item $\mathbf{P}_e(\mathbf{X},t)$: \emph{first Piola--Kirchhoff elastic stress}
\end{itemize}

\medskip

The first Piola--Kirchhoff stress is defined by
\[
\mathbf{P}_e(\mathbf{X},t)
=
\frac{\partial W_e}{\partial \mathbf{F}}(\mathbf{F}(\mathbf{X},t)),
\qquad
\mathbf{F} = \frac{\partial \boldsymbol{\chi}}{\partial \mathbf{X}}
\]

where $W_e$ is the elastic strain-energy density.

\medskip

\textbf{Key point:}  
$\mathbf{P}_e$ is a stress measured \emph{per unit reference area}, not per unit physical area.

\end{frame}

\begin{frame}{Strong Form: From Lagrangian Stress to Eulerian Force}

The Lagrangian elastic force density is given by the material divergence:
\[
\mathbf{F}_e(\mathbf{X},t)
=
\nabla_{\mathbf{X}} \cdot \mathbf{P}_e(\mathbf{X},t)
\]

This quantity has units of force per unit \emph{reference volume}.

\medskip

The corresponding Eulerian force density is obtained via force spreading:
\[
\mathbf{f}_e(\mathbf{x},t)
=
\int_{\mathcal{U}}
\mathbf{F}_e(\mathbf{X},t)\,
\delta(\mathbf{x}-\boldsymbol{\chi}(\mathbf{X},t))\,d\mathbf{X}
\]

\medskip

Interpretation:
\begin{itemize}
    \item Forces are computed in reference coordinates
    \item The delta function deposits those forces at their current physical locations
\end{itemize}

\end{frame}

%=======================================================
\section{Why No Jacobian Appears in the IB Force Spreading}
%=======================================================

\begin{frame}{Why No Jacobian Appears}
\small

Consider the classical IB force-spreading formula (Peskin):
\[
\mathbf{f}_e(\mathbf{x},t)
=
\int_{\mathcal{U}}
\mathbf{F}_e(\mathbf{X},t)\,
\delta\!\big(\mathbf{x}-\boldsymbol{\chi}(\mathbf{X},t)\big)\,d\mathbf{X}.
\]

\medskip

\textbf{What is important:}
\begin{itemize}
    \item The integral is performed entirely over the \emph{reference (material) domain} $\mathcal{U}$
    \item $\mathbf{F}_e(\mathbf{X},t)$ is defined as a force density \emph{per unit reference measure} $d\mathbf{X}$
    \item No change of variables from $\mathbf{X}$ to $\mathbf{x}$
    \item No rewriting of $d\mathbf{X}$ in terms of $d\mathbf{x}$
\end{itemize}

\[
d\mathbf{X} \neq J^{-1}\,d\mathbf{x},
\qquad
J = \det\!\left(\frac{\partial\boldsymbol{\chi}}{\partial\mathbf{X}}\right),
\]
because the integral is \emph{not} expressed in Eulerian coordinates.

\medskip

\textbf{Peskin (1996):}
The delta-function formulation ``avoids the introduction of the Jacobian'' by defining force spreading as a distributional mapping from material to physical space.

\end{frame}

%=======================================================

\begin{frame}{Delta Function as a Push-Forward (Not a Change of Variables)}
\small

The Dirac delta function
\[
\delta\!\big(\mathbf{x}-\boldsymbol{\chi}(\mathbf{X},t)\big)
\]
acts as a \emph{distributional push-forward} from material space to physical space.

\medskip

Formally, for any smooth test function $\boldsymbol{\varphi}(\mathbf{x})$,
\[
\int_{\Omega}\boldsymbol{\varphi}(\mathbf{x})\,\mathbf{f}_e(\mathbf{x},t)\,d\mathbf{x}
=
\int_{\mathcal{U}}
\boldsymbol{\varphi}\!\big(\boldsymbol{\chi}(\mathbf{X},t)\big)\,
\mathbf{F}_e(\mathbf{X},t)\,d\mathbf{X}.
\]

\medskip

\textbf{Key point (Peskin):}
This identity defines $\mathbf{f}_e$ as the Eulerian force density whose action on test
functions equals the action of the Lagrangian force under the mapping $\boldsymbol{\chi}$.

\medskip

\textbf{Therefore:}
\begin{itemize}
    \item Geometry enters through $\boldsymbol{\chi}$ and the delta distribution
    \item No Jacobian is required because no coordinate transformation is performed
\end{itemize}

\end{frame}

%=======================================================

\begin{frame}{When Would a Jacobian (or Metric Factor) Appear?}
\small

The classical IB force spreading
\[
\mathbf{f}(\mathbf{x},t)
=
\int_{\mathcal{U}}
\mathbf{F}_{\mathrm{ref}}(\mathbf{X},t)\,
\delta\!\big(\mathbf{x}-\boldsymbol{\chi}(\mathbf{X},t)\big)\,d\mathbf{X}
\]
assumes that $\mathbf{F}_{\mathrm{ref}}$ is a force density per unit
\emph{reference} measure.

\medskip

A Jacobian or metric factor appears when the force density is instead defined per unit
\emph{current (deformed) physical measure}.

\medskip

\textbf{Examples:}
\begin{itemize}
    \item \textbf{Volumetric solid ($k=d$):}
    \[
    d\mathbf{x} = J\,d\mathbf{X},
    \qquad
    J=\det\!\left(\frac{\partial\boldsymbol{\chi}}{\partial\mathbf{X}}\right)
    \]
    \item \textbf{Surface in 3D ($k=2$):}
    \[
    da = J_s\,dA \quad \text{(surface Jacobian / metric)}
    \]
    \item \textbf{Fiber in 2D or 3D ($k=1$):}
    \[
    ds = \lambda\,dS,
    \qquad
    \lambda = \left|\frac{\partial\boldsymbol{\chi}}{\partial S}\right|
    \]
\end{itemize}

\medskip

\textbf{Key point:}
“No Jacobian” is correct \emph{only} when forces are defined per unit reference measure.
Metric factors enter when the model is posed in current measures.

\end{frame}

%=======================================================

\begin{frame}{Reference vs.\ Current Force Densities: Same Physical Force}
\small

Let
\begin{itemize}
    \item $\mathbf{F}_{\mathrm{ref}}(\mathbf{X},t)$ be force per unit reference measure $d\mathbf{X}$
    \item $\mathbf{F}_{\mathrm{cur}}(\mathbf{X},t)$ be force per unit current physical measure
\end{itemize}

\medskip

To represent the \emph{same total physical force}, we require
\[
\mathbf{F}_{\mathrm{ref}}(\mathbf{X},t)\,d\mathbf{X}
=
\mathbf{F}_{\mathrm{cur}}(\mathbf{X},t)\,d\mathbf{x}
\quad
(\text{or } da,\, ds).
\]

\medskip

\textbf{Volumetric solid:}
\[
d\mathbf{x}=J\,d\mathbf{X}
\quad\Longrightarrow\quad
\mathbf{F}_{\mathrm{ref}} = J\,\mathbf{F}_{\mathrm{cur}}.
\]

\medskip

Thus the same Eulerian force density may be written as
\[
\mathbf{f}(\mathbf{x},t)
=
\int_{\mathcal{U}}
\mathbf{F}_{\mathrm{ref}}(\mathbf{X},t)\,
\delta(\mathbf{x}-\boldsymbol{\chi})\,d\mathbf{X}
=
\int_{\mathcal{U}}
\mathbf{F}_{\mathrm{cur}}(\mathbf{X},t)\,
\delta(\mathbf{x}-\boldsymbol{\chi})\,J\,d\mathbf{X}.
\]

\medskip

\textbf{Conclusion:}
This is a modeling choice, not a contradiction of the IB formulation.

\end{frame}

%=======================================================

% Requires in preamble:
% \usepackage{amsmath,amssymb}
% \usepackage{tikz}
% \usetikzlibrary{arrows.meta,positioning,calc}

% ------------------------------------------------------------

% In your preamble:
% \usepackage{amsmath,amssymb}
% \usepackage{tikz}

% ------------------------------------------------------------
\begin{frame}{Force spreading: the definition (IB interaction equation)}
\footnotesize

Let $\boldsymbol{\chi}(\mathbf{X},t)$ map Lagrangian labels $\mathbf{X}\in\mathcal{U}$ to physical space
$\mathbf{x}\in\Omega$, and let $\mathbf{F}(\mathbf{X},t)$ be a Lagrangian force density.

\medskip
\textbf{Definition (standard in IBM / IB methods):}
\[
\boxed{
\mathbf{f}(\mathbf{x},t)
=
\int_{\mathcal{U}}
\mathbf{F}(\mathbf{X},t)\,
\delta\!\big(\mathbf{x}-\boldsymbol{\chi}(\mathbf{X},t)\big)\,d\mathbf{X}
}
\qquad (\mathbf{x}\in\Omega).
\]

\medskip
Here $\delta(\mathbf{x})=\prod_{i=1}^{d}\delta(x_i)$ is the $d$-dimensional Dirac delta ($d=2$ or $3$).

\medskip
\textbf{Goal:} show this definition means
\emph{“total Eulerian force on any region equals the total Lagrangian force carried by points in that region.”}
\end{frame}

% ------------------------------------------------------------
\begin{frame}{Force spreading: test with a region $\Gamma$ (Step 1--2)}
\footnotesize

Pick any Eulerian region $\Gamma\subset\Omega$ and integrate over it.

\medskip
\textbf{Step 1: integrate the Eulerian force density over $\Gamma$:}
\[
\int_{\Gamma}\mathbf{f}(\mathbf{x},t)\,d\mathbf{x}.
\]

\medskip
\textbf{Step 2: substitute the definition of $\mathbf{f}$:}
\[
\int_{\Gamma}\mathbf{f}(\mathbf{x},t)\,d\mathbf{x}
=
\int_{\Gamma}\int_{\mathcal{U}}
\mathbf{F}(\mathbf{X},t)\,
\delta\!\big(\mathbf{x}-\boldsymbol{\chi}(\mathbf{X},t)\big)\,d\mathbf{X}\,d\mathbf{x}.
\]

\medskip
(Think of this as: sum up contributions from all material points, but each point only contributes where the delta “hits”.)
\end{frame}

% ------------------------------------------------------------
\begin{frame}{Force spreading: swap integrals + use the delta (Step 3--4)}
\footnotesize

\textbf{Step 3: swap the order of integration:} (see next slides for justification)
\[
\int_{\Gamma}\int_{\mathcal{U}}(\cdot)\,d\mathbf{X}\,d\mathbf{x}
=
\int_{\mathcal{U}}
\mathbf{F}(\mathbf{X},t)\left(
\int_{\Gamma}\delta\!\big(\mathbf{x}-\boldsymbol{\chi}(\mathbf{X},t)\big)\,d\mathbf{x}
\right)\,d\mathbf{X}.
\]

\medskip
\textbf{Step 4: apply the defining property of the delta distribution:}
\[
\int_{\Gamma}\delta\!\big(\mathbf{x}-\boldsymbol{\chi}(\mathbf{X},t)\big)\,d\mathbf{x}
=
\begin{cases}
1, & \boldsymbol{\chi}(\mathbf{X},t)\in \Gamma,\\
0, & \boldsymbol{\chi}(\mathbf{X},t)\notin \Gamma.
\end{cases}
\]

\medskip
So only those material points whose \emph{current position} lies in $\Gamma$ contribute to the integral.
\end{frame}

% ------------------------------------------------------------
\begin{frame}{Force spreading: the resulting identity (the punchline)}
\footnotesize

Combining the previous steps yields:
\[
\boxed{
\int_{\Gamma}\mathbf{f}(\mathbf{x},t)\,d\mathbf{x}
=
\int_{\{\mathbf{X}\in\mathcal{U}:\,\boldsymbol{\chi}(\mathbf{X},t)\in\Gamma\}}
\mathbf{F}(\mathbf{X},t)\,d\mathbf{X}.
}
\]

\medskip
\textbf{Interpretation (Peskin):}
The Eulerian force density $\mathbf{f}$ is defined so that the \emph{total force applied to any fluid region $\Gamma$}
equals the \emph{total Lagrangian force carried by material points occupying that region}.
\end{frame}

% ------------------------------------------------------------


% ------------------------------------------------------------
\begin{frame}{One-slide summary: what you should remember}
\footnotesize

\textbf{1) Spread Lagrangian force to the Eulerian grid:}
\[
\boxed{
\mathbf{f}(\mathbf{x},t)=\int_{\mathcal{U}}\mathbf{F}(\mathbf{X},t)\,
\delta\!\big(\mathbf{x}-\boldsymbol{\chi}(\mathbf{X},t)\big)\,d\mathbf{X}.
}
\]

\medskip
\textbf{2) Meaning (test with any region $\Gamma\subset\Omega$):}
\[
\boxed{
\int_{\Gamma}\mathbf{f}(\mathbf{x},t)\,d\mathbf{x}
=
\int_{\chi^{-1}(\Gamma)}\mathbf{F}(\mathbf{X},t)\,d\mathbf{X}.
}
\]
So the \emph{total force on a fluid region} equals the \emph{total force carried by material points in that region}.

\medskip
\textbf{3) In computations:}
replace $\delta$ by a smooth, compactly-supported $\delta_h$ (regularized delta) so spreading/interpolation
happen over a few nearby grid cells.
\end{frame}

% ------------------------------------------------------------
\begin{frame}{Can we swap the order of integration? (The subtlety)}
\footnotesize

We start from the IB force-spreading identity
\[
\mathbf{f}(\mathbf{x},t)=\int_{\mathcal U}\mathbf{F}(\mathbf{X},t)\,
\delta\!\big(\mathbf{x}-\boldsymbol{\chi}(\mathbf{X},t)\big)\,d\mathbf{X}.
\]

\medskip
If we integrate over a region $\Gamma\subset\Omega$, we are tempted to write
\[
\int_{\Gamma}\mathbf{f}(\mathbf{x},t)\,d\mathbf{x}
=
\int_{\Gamma}\int_{\mathcal U}\mathbf{F}(\mathbf{X},t)\,
\delta\!\big(\mathbf{x}-\boldsymbol{\chi}(\mathbf{X},t)\big)\,d\mathbf{X}\,d\mathbf{x}
\overset{?}{=}
\int_{\mathcal U}\int_{\Gamma}(\cdots)\,d\mathbf{x}\,d\mathbf{X}.
\]

\medskip
\textbf{Key point:} the Dirac delta is \emph{not an $L^1$ function}. So classical Fubini/Tonelli
does not directly apply to the double integral with $\delta$.

\medskip
\textbf{Two standard ways forward:}
\begin{enumerate}\itemsep2pt
\item Interpret the formula in the \textbf{weak (distributional) sense}. No ``swapping'' is needed.
\item Replace $\delta$ by a \textbf{regularized kernel $\delta_h$} (what we do numerically),
      and then Fubini is justified by absolute integrability.
\end{enumerate}

\end{frame}

% ------------------------------------------------------------
\begin{frame}{Justification 1: weak/distribution interpretation}
\footnotesize

\textbf{Define $\mathbf{f}$ as a distribution.}
Let $\boldsymbol{\varphi}(\mathbf{x})\in C_c^\infty(\Omega)$ be a smooth test function (vector-valued).
Then define $\mathbf{f}$ by its action on $\boldsymbol{\varphi}$:
\[
\int_{\Omega}\mathbf{f}(\mathbf{x},t)\cdot \boldsymbol{\varphi}(\mathbf{x})\,d\mathbf{x}
:=
\int_{\mathcal U}\mathbf{F}(\mathbf{X},t)\cdot
\boldsymbol{\varphi}\!\big(\boldsymbol{\chi}(\mathbf{X},t)\big)\,d\mathbf{X}.
\]

\medskip
\textbf{This is exactly what the delta notation means:}
\[
\int_{\Omega}\left(\int_{\mathcal U}\mathbf{F}(\mathbf{X},t)\,
\delta(\mathbf{x}-\boldsymbol{\chi}(\mathbf{X},t))\,d\mathbf{X}\right)\cdot
\boldsymbol{\varphi}(\mathbf{x})\,d\mathbf{x}
=
\int_{\mathcal U}\mathbf{F}(\mathbf{X},t)\cdot \boldsymbol{\varphi}(\boldsymbol{\chi}(\mathbf{X},t))\,d\mathbf{X}.
\]

\medskip
\textbf{Region version (idea):} choose a sequence $\boldsymbol{\varphi}_\epsilon$ that approximates
the indicator of $\Gamma$ (smoothed characteristic). Passing to the limit yields
\[
\int_{\Gamma}\mathbf{f}(\mathbf{x},t)\,d\mathbf{x}
=
\int_{\{\mathbf{X}:\,\boldsymbol{\chi}(\mathbf{X},t)\in \Gamma\}}\mathbf{F}(\mathbf{X},t)\,d\mathbf{X}.
\]

\medskip
\textbf{Bottom line:} with a true Dirac delta, the identity is understood in a weak sense,
not as an $L^1$ double integral where we invoke Fubini.

\end{frame}



%=======================================================

\begin{frame}{Intuition: Force Spreading as a Push-Forward}
\small

The force-spreading integral
\[
\mathbf{f}(\mathbf{x},t)
=
\int_{\mathcal{U}}
\mathbf{F}(\mathbf{X},t)\,
\delta\!\big(\mathbf{x}-\boldsymbol{\chi}(\mathbf{X},t)\big)\,d\mathbf{X}
\]
is best understood as a \emph{push-forward of forces} from material space to physical space.

\medskip

\textbf{Interpretation:}
\begin{itemize}
    \item Forces are computed on the structure in reference coordinates
    \item $\boldsymbol{\chi}(\mathbf{X},t)$ determines where those forces act in space
    \item The delta distribution transfers those forces to the fluid
\end{itemize}

\medskip

\textbf{Crucial point (Peskin):}
No Jacobian appears because this is \emph{not} a change of variables in an integral.
All integration is performed in reference coordinates; geometry enters only through
$\boldsymbol{\chi}$ and the delta distribution.

\medskip

\textit{Total force on a fluid region equals the total force carried by structure points inside that region.}

\end{frame}



%=======================================================
\begin{frame}{Units Check: 1D Structure in a 2D Fluid}

Consider a 1D immersed structure (fiber or membrane)
in a 2D incompressible fluid.

\medskip

Force spreading:
\[
\mathbf{f}(\mathbf{x})
=
\int
\mathbf{F}(s)\,
\delta(\mathbf{x}-\mathbf{X}(s))\,ds
\]

\medskip

Units:
\[
[\mathbf{F}] = \frac{\text{N}}{\text{m}},
\qquad
[\delta] = \frac{1}{\text{m}^2},
\qquad
[ds] = \text{m}
\]

\medskip

Thus:
\[
[\mathbf{f}]
=
\frac{\text{N}}{\text{m}}
\cdot
\frac{1}{\text{m}^2}
\cdot
\text{m}
=
\frac{\text{N}}{\text{m}^2}
\]

\medskip

which is the correct Eulerian force per unit area.

\end{frame}

\begin{frame}{Units Check: General Case}
\small
Let:
\begin{itemize}
    \item $d$ = Eulerian spatial dimension ($d=2$ or $3$)
    \item $k$ = Lagrangian dimension of the structure ($k \le d$)
    \item $\mathbf{F}(\mathbf{X})$ = Lagrangian force density
\end{itemize}

\medskip

Force spreading:
\[
\mathbf{f}(\mathbf{x})
=
\int_{\mathcal{U}}
\mathbf{F}(\mathbf{X})\,
\delta(\mathbf{x}-\boldsymbol{\chi}(\mathbf{X}))\,d\mathbf{X}
\]

\medskip

Units:
\[
[\mathbf{F}] = \frac{\text{N}}{\text{m}^k},
\qquad
[\delta] = \frac{1}{\text{m}^d},
\qquad
[d\mathbf{X}] = \text{m}^k
\]

\medskip

Thus:
\[
[\mathbf{f}]
=
\frac{\text{N}}{\text{m}^k}
\cdot
\frac{1}{\text{m}^d}
\cdot
\text{m}^k
=
\frac{\text{N}}{\text{m}^d}
\]

\medskip

which is the correct Eulerian force per unit $d$-dimensional volume
(area in 2D, volume in 3D). $k=1$ (filament), $k=2$ (surface), $k=3$ (solid).

\end{frame}

%=======================================================
\section{Intuition for Force Spreading}
\begin{frame}{Intuition: What Does Force Spreading Mean?}

The force-spreading transform
\[
\mathbf{f}(\mathbf{x},t)
=
\int_{\mathcal{U}}
\mathbf{F}(\mathbf{X},t)\,
\delta(\mathbf{x}-\boldsymbol{\chi}(\mathbf{X},t))\,d\mathbf{X}
\]
has a simple physical interpretation.

\medskip

\textbf{Think of it as:}
\begin{itemize}
    \item Each material point $\mathbf{X}$ carries a force $\mathbf{F}(\mathbf{X},t)\,d\mathbf{X}$
    \item That force is applied to the fluid at the point
    \[
    \mathbf{x}=\boldsymbol{\chi}(\mathbf{X},t)
    \]
    \item The delta function places the force at the correct physical location
\end{itemize}

\medskip

In words:
\[
\text{``Compute forces on the structure, then deposit them into the fluid.''}
\]

\end{frame}

\begin{frame}{Intuition: A Distribution, Not a Point Force}

The Eulerian force density $\mathbf{f}(\mathbf{x},t)$ is \emph{not} a regular function.

\medskip

It is a \textbf{distribution} supported on the immersed structure:
\[
\mathbf{f}(\mathbf{x},t)
=
\int_{\mathcal{U}}
\mathbf{F}(\mathbf{X},t)\,
\delta(\mathbf{x}-\boldsymbol{\chi}(\mathbf{X},t))\,d\mathbf{X}
\]

\medskip

Interpretation:
\begin{itemize}
    \item The force is concentrated on a lower-dimensional set
    \item In the continuum limit, this produces a $\delta$-layer of force
    \item Physically: a thin membrane, fiber, or surface pushing on the fluid
\end{itemize}

\medskip

This is why:
\begin{itemize}
    \item Pressure and stress can be discontinuous across the structure
    \item Velocity remains continuous (viscous fluid)
\end{itemize}

\end{frame}


\begin{frame}{Intuition: Regularized Delta Functions}

In computation, the Dirac delta is replaced by a smooth kernel $\delta_h$:
\[
\mathbf{f}(\mathbf{x})
\approx
\int_{\mathcal{U}}
\mathbf{F}(\mathbf{X})\,
\delta_h(\mathbf{x}-\boldsymbol{\chi}(\mathbf{X}))\,d\mathbf{X}
\]

\medskip

Interpretation:
\begin{itemize}
    \item A single Lagrangian force is \emph{shared} among nearby grid points
    \item The force is spread over a region of width $\mathcal{O}(h)$
    \item This avoids singular forces on a Cartesian grid
\end{itemize}

\medskip

Key idea:
\[
\text{Force spreading} \;\;=\;\; \text{local averaging consistent with conservation}
\]

\medskip

The same kernel is used for:
\begin{itemize}
    \item Force spreading (Lagrangian $\rightarrow$ Eulerian)
    \item Velocity interpolation (Eulerian $\rightarrow$ Lagrangian)
\end{itemize}

\end{frame}


%=======================================================
\section{Regularized Delta Functions}
%=======================================================

\begin{frame}{From the Dirac Delta to a Discrete Kernel}

In the continuum immersed boundary method, interaction is mediated by
the Dirac delta function:
\[
\mathbf{f}(\mathbf{x},t)
=
\int_{\mathcal{U}}
\mathbf{F}(\mathbf{X},t)\,
\delta(\mathbf{x}-\boldsymbol{\chi}(\mathbf{X},t))\,d\mathbf{X}.
\]

\medskip

In numerical computations:
\begin{itemize}
    \item The Eulerian domain is discretized on a Cartesian grid with spacing $h$
    \item The Lagrangian structure is represented by discrete points $\mathbf{X}_k$
    \item The singular delta distribution must be replaced by a smooth approximation
\end{itemize}

\medskip

This leads to the use of a \textbf{regularized (discrete) delta function}
$\delta_h$.

\end{frame}

\begin{frame}{Discrete Interaction via a Regularized Delta Function}

Using a regularized delta function $\delta_h$, interaction is approximated as:

\medskip

\textbf{Velocity interpolation (Eulerian $\rightarrow$ Lagrangian):}
\[
\mathbf{U}_k
=
\sum_{\mathbf{i}}
\mathbf{u}_{\mathbf{i}}\,
\delta_h(\mathbf{x}_{\mathbf{i}}-\mathbf{X}_k)\,h^d
\]

\medskip

\textbf{Force spreading (Lagrangian $\rightarrow$ Eulerian):}
\[
\mathbf{f}_{\mathbf{i}}
=
\sum_k
\mathbf{F}_k\,
\delta_h(\mathbf{x}_{\mathbf{i}}-\mathbf{X}_k)\,\Delta V_k
\]

\medskip

where:
\begin{itemize}
    \item $h$ is the Eulerian grid spacing
    \item $\Delta V_k$ is the Lagrangian quadrature weight
    \item $d=2$ or $3$ is the Eulerian spatial dimension
\end{itemize}

\end{frame}

\begin{frame}{Structure of the Discrete Delta Function}

In practice, the discrete delta function is built as a tensor product of
1D kernel functions:
\[
\delta_h(\mathbf{x})
=
\frac{1}{h^d}
\prod_{\ell=1}^{d}
\phi\!\left(\frac{x_\ell}{h}\right).
\]

\medskip

Here:
\begin{itemize}
    \item $\phi(r)$ is a smooth, compactly supported 1D kernel
    \item The full kernel has support over a small number of grid cells
    \item The same kernel is used in each coordinate direction
\end{itemize}

\medskip

This construction preserves isotropy on a Cartesian grid and simplifies implementation.

\end{frame}

\begin{frame}{Why Special Regularized Delta Functions Are Needed}

In the immersed boundary method, the Dirac delta distribution
must be replaced by a smooth approximation $\delta_h$.

\medskip

This approximation must:
\begin{itemize}
    \item faithfully represent interaction between Eulerian and Lagrangian grids
    \item preserve basic physical invariants (force, torque)
    \item be compatible with a Cartesian grid
    \item remain numerically stable
\end{itemize}

\medskip

Peskin (1996) proposed a set of \textbf{design criteria} for $\delta_h$
that balance accuracy, smoothness, and efficiency.

\end{frame}

\begin{frame}{Property 1: Compact Support (Locality)}

The discrete delta function must have compact support.

\medskip

In one dimension:
\[
\phi(r) = 0 \quad \text{for } |r| \geq 2.
\]

\medskip

In $d$ dimensions:
\[
\delta_h(\mathbf{x})
=
\frac{1}{h^d}
\prod_{\ell=1}^{d}
\phi\!\left(\frac{x_\ell}{h}\right),
\]
so each Lagrangian point interacts with only a small stencil of grid points.

\medskip

\textbf{Why this matters:}
\begin{itemize}
    \item limits computational cost
    \item prevents nonphysical long-range interactions
    \item yields sparse coupling between grids
\end{itemize}

\end{frame}

\begin{frame}{Property 2: Zeroth Moment (Force Conservation)}

The discrete delta must satisfy a discrete partition of unity:
\[
\sum_{i}
\phi(r-i)
=
1
\qquad \text{for all } r.
\]

\medskip

Equivalently,
\[
\sum_i \delta_h(r-ih)\,h = 1.
\]

\medskip

\textbf{Physical meaning:}
\begin{itemize}
    \item total force applied to the fluid equals total force generated by the structure
    \item no artificial loss or gain of force during spreading
\end{itemize}

\medskip

This condition is essential for global momentum conservation.

\end{frame}

\begin{frame}{Property 3: First Moment (No Spurious Torque)}

The discrete delta must have zero first moment:
\[
\sum_i (r-i)\,\phi(r-i) = 0.
\]

\medskip

Equivalently,
\[
\sum_i (r-ih)\,\delta_h(r-ih)\,h = 0.
\]

\medskip

\textbf{Physical meaning:}
\begin{itemize}
    \item the force distribution is centered correctly
    \item no artificial torque is introduced by spreading
\end{itemize}

\medskip

This is especially important for:
\begin{itemize}
    \item rotating structures
    \item bending-dominated elastic models
\end{itemize}

\end{frame}

\begin{frame}{Property 4: Symmetry and Translational Invariance}

Peskin requires the 1D kernel to be even:
\[
\phi(r) = \phi(-r).
\]

\medskip

This ensures:
\begin{itemize}
    \item symmetry of force spreading
    \item approximate translational invariance on a Cartesian grid
\end{itemize}

\medskip

\textbf{Interpretation:}
\begin{itemize}
    \item numerical results should not depend strongly on the exact
          position of the structure relative to grid points
    \item avoids grid-induced directional bias
\end{itemize}

\end{frame}

\begin{frame}{Property 5: Smoothness and Energy Consistency}

Peskin also imposes a \textbf{sum-of-squares condition}:
\[
\sum_i \phi^2(r-i) = C,
\]
where $C$ is a constant independent of $r$.

\medskip

\textbf{Consequences:}
\begin{itemize}
    \item improves numerical stability
    \item supports consistent coupling between interpolation and spreading
    \item reduces grid-scale oscillations
\end{itemize}

\medskip

Additionally, $\phi$ is chosen to be as smooth as possible
(subject to compact support).

\medskip

\textbf{Tradeoff:}  
More smoothness $\Rightarrow$ wider support $\Rightarrow$ higher cost.

\end{frame}

\begin{frame}{Why the Same Kernel is Used Both Ways}
\small

In the discrete method, spreading and interpolation should be compatible so that
the work done on the fluid matches the work extracted from the structure.

\medskip
Discrete fluid power input:
\[
\sum_{\mathbf{i}} \mathbf{u}(\mathbf{x}_{\mathbf{i}})\cdot \mathbf{f}(\mathbf{x}_{\mathbf{i}})\,h^d.
\]

Using force spreading
\[
\mathbf{f}(\mathbf{x}_{\mathbf{i}})
=
\sum_k \mathbf{F}_k\,\delta_h(\mathbf{x}_{\mathbf{i}}-\mathbf{X}_k)\,\Delta V_k,
\]
we obtain
\[
\sum_{\mathbf{i}} \mathbf{u}(\mathbf{x}_{\mathbf{i}})\cdot \mathbf{f}(\mathbf{x}_{\mathbf{i}})\,h^d
=
\sum_k \mathbf{F}_k \cdot
\left(
\sum_{\mathbf{i}}\mathbf{u}(\mathbf{x}_{\mathbf{i}})\,\delta_h(\mathbf{x}_{\mathbf{i}}-\mathbf{X}_k)\,h^d
\right)\Delta V_k.
\]

\medskip
The term in parentheses is exactly the \emph{interpolated velocity} $\mathbf{U}_k$.
So we get:
\[
\sum_{\mathbf{i}} \mathbf{u}\cdot \mathbf{f}\,h^d
=
\sum_k \mathbf{F}_k\cdot \mathbf{U}_k\,\Delta V_k.
\]

\medskip
This motivates using the same kernel and carefully chosen kernel properties.

\end{frame}


\begin{frame}{Summary: Peskin’s Delta Function Design Principles}

Peskin’s regularized delta functions are:
\begin{itemize}
    \item compactly supported
    \item force-conserving (zeroth moment)
    \item torque-free (first moment)
    \item symmetric and approximately translationally invariant
    \item smooth and numerically stable
\end{itemize}

\medskip

The resulting kernels are not unique, but the 4-point kernel
provides an effective balance between:
\[
\text{accuracy} \;\;+\;\; \text{stability} \;\;+\;\; \text{efficiency}.
\]

\end{frame}

\begin{frame}{Discrete Delta Functions Used in Practice}

In the immersed boundary method, the discrete delta function
is constructed from a 1D kernel $\phi(r)$ as:
\[
\delta_h(\mathbf{x})
=
\frac{1}{h^d}
\prod_{\ell=1}^{d}
\phi\!\left(\frac{x_\ell}{h}\right).
\]

\medskip

Different choices of $\phi(r)$ satisfy Peskin’s design criteria
to varying degrees.

\medskip

We now present several commonly used kernels and discuss their properties.

\end{frame}

\begin{frame}{Peskin’s 4-Point Delta Function}

The most widely used kernel in immersed boundary methods
(Peskin, 1996) has support over four grid points.

\medskip

The 1D kernel is:
\[
\phi(r)
=
\begin{cases}
\frac{1}{8}
\left(
3 - 2|r| + \sqrt{1 + 4|r| - 4r^2}
\right),
& 0 \le |r| < 1, \\[6pt]
\frac{1}{8}
\left(
5 - 2|r| - \sqrt{-7 + 12|r| - 4r^2}
\right),
& 1 \le |r| < 2, \\[6pt]
0,
& |r| \ge 2.
\end{cases}
\]

\medskip

\textbf{Properties:}
\begin{itemize}
    \item Compact support: $|r| < 2$
    \item Satisfies zeroth and first moment conditions
    \item Satisfies Peskin’s sum-of-squares condition
    \item Good balance of smoothness and locality
\end{itemize}

\end{frame}

\begin{frame}{Peskin / Roma 3-Point Delta Function}

A smaller-support alternative is the 3-point kernel
(Roma, Peskin, Berger).

\medskip

The 1D kernel is:
\[
\phi(r)
=
\begin{cases}
\frac{1}{3}
\left(
1 + \sqrt{1 - 3r^2}
\right),
& |r| < \frac{1}{2}, \\[6pt]
\frac{1}{6}
\left(
5 - 3|r| - \sqrt{-2 + 6|r| - 3r^2}
\right),
& \frac{1}{2} \le |r| < \frac{3}{2}, \\[6pt]
0,
& |r| \ge \frac{3}{2}.
\end{cases}
\]

\medskip

\textbf{Properties:}
\begin{itemize}
    \item Smaller support than the 4-point kernel
    \item Slightly less smooth
    \item Lower computational cost
\end{itemize}

\end{frame}

\begin{frame}{Cosine Delta Function}

A simpler alternative kernel uses a cosine profile.

\medskip

The 1D kernel is:
\[
\phi(r)
=
\begin{cases}
\frac{1}{4}
\left(
1 + \cos\!\left(\frac{\pi r}{2}\right)
\right),
& |r| \le 2, \\[6pt]
0,
& |r| > 2.
\end{cases}
\]

\medskip

\textbf{Properties:}
\begin{itemize}
    \item Compact support over $|r|\le 2$
    \item $C^1$ continuous
    \item Easy to implement
    \item Does \emph{not} satisfy Peskin’s sum-of-squares condition exactly
\end{itemize}

\medskip

Often used for:
\begin{itemize}
    \item exploratory simulations
    \item pedagogical implementations
\end{itemize}

\end{frame}

\begin{frame}{Higher-Order Delta Functions}

Higher-order kernels with wider support have also been proposed.

\medskip

Examples:
\begin{itemize}
    \item 6-point Peskin kernel
    \item B-spline-based kernels
\end{itemize}

\medskip

\textbf{Advantages:}
\begin{itemize}
    \item Increased smoothness
    \item Improved translational and rotational invariance
\end{itemize}

\medskip

\textbf{Disadvantages:}
\begin{itemize}
    \item Larger computational stencil
    \item Higher cost per time step
\end{itemize}

\medskip

Choice of kernel depends on the accuracy–cost tradeoff
appropriate for the problem.

\end{frame}

\begin{frame}{Which Delta Function Should You Use?}

\begin{itemize}
    \item \textbf{Peskin 4-point}: default choice in IBAMR and IB2d
    \item \textbf{3-point}: faster, but less smooth
    \item \textbf{Cosine}: simple, intuitive, but less theoretically robust
    \item \textbf{Higher-order}: best accuracy, higher cost
\end{itemize}

\medskip

\textbf{Rule of thumb:}
\begin{itemize}
    \item Start with the 4-point kernel
    \item Change kernels only if you have a clear reason
\end{itemize}

\end{frame}




\begin{frame}{Physical Interpretation of $\delta_h$}

The regularized delta function does \emph{not} represent a point force.

\medskip

Instead, it:
\begin{itemize}
    \item distributes a Lagrangian force smoothly over nearby grid points
    \item approximates the action of a lower-dimensional structure on the fluid
    \item avoids singular forces on a Cartesian grid
\end{itemize}

\medskip

As $h \rightarrow 0$:
\[
\delta_h \;\longrightarrow\; \delta
\]

and the discrete method converges to the continuum immersed boundary formulation.

\end{frame}

%=======================================================
\begin{frame}
\centering
\Large Questions?
\end{frame}

\end{document}
